\subsection{Lyapunov quantification of the closed model}
\begingroup\color{red}
We quantify sensitivity to initial histories in the underlying DDE.
The calculation concerns the closed model with fixed coefficients,
not the neural network or the observed activity series.
Time uses the normalized units of the empirical implementation.
For $\tau>0$, a perturbation satisfies
\begin{equation}
\dot v(t)=A(t)v(t)+B(t)v(t-\tau),
\label{eq:lyapunov_variation}
\end{equation}
where $A(t)=\partial_1 f$ and $B(t)=\partial_2 f$ are evaluated
along the DDE solution.
The seven human components of the initial perturbation sum to zero
at every history time.
For $v_t(s)=v(t+s)$, use
\[
\|v_t\|_\tau^2
=\frac{1}{\tau}\int_{-\tau}^{0}\|v(t+s)\|_2^2\,ds.
\]
Let $\mathcal T$ denote this space of perturbation histories.
The maximal upper growth exponent is
\[
\overline\lambda_{\max}(X)
=\sup_{v_0\in\mathcal T\setminus\{0\}}
\limsup_{T\to\infty}
\frac{1}{T}\log\frac{\|v_T\|_\tau}{\|v_0\|_\tau}.
\]
This definition does not assume convergence of a finite duration estimate.
We estimate the dominant rate by propagating several initial directions.
At each normalization time $t_k$, we divide the complete perturbation
history by its norm $g_k$~\cite{ansmann2018lyapunov}.
For $t_j=b$ and $t_K=T$, the estimate is
\begin{equation}
\widehat\lambda_{b,T}
=\frac{1}{T-b}\sum_{k=j+1}^{K}\log g_k.
\label{eq:lyapunov_growth_identity}
\end{equation}
Linearity makes this sum equal to the logarithmic growth of the
corresponding history without normalization.
It is a directional estimate, not a maximum over all histories
at finite $T$.

\begin{theorem}[Extinction and Lyapunov growth in the closed model]
\label{thm:lyapunov_quantification}
Consider the closed model with $N=1$ and $\Lambda=d_h=0$.
Use the fixed empirical coefficients
\[
\begin{gathered}
\gamma_E=0.2,\quad \mu=\alpha=0.25,\quad q_U=0.3,\\
\xi_E=0.8,\quad \sigma_E=0.4,\quad
\kappa_1=0.4,\quad \kappa_2=0.25,\\
\eta_1=\eta_2=0.6,\quad v_1=v_2=0.2,\quad
d_{M_1}=d_{M_2}=1,\\
\delta_1=0.5,\quad \delta_2=0.25,\quad
\varepsilon=0.55,\quad \rho=5,\quad
\varpi_1=\varpi_2=1.
\end{gathered}
\]
All remaining coefficients are finite and nonnegative, with $\chi>0$.
Initial histories are nonnegative and satisfy the cohort compatibility
condition in Theorem~III.1.
Every solution converges to an equilibrium
\begin{equation}
X^*(s)=
(5s/13,8s/13,0,0,0,0,1-s,0.2,0.2),
\qquad 0\le s\le1.
\label{eq:closed_equilibrium_family}
\end{equation}
Thus, this closed model has no sustained chaotic attractor.

At $X^*(s)$, set
\[
\begin{gathered}
\psi_s=\frac{47s}{65},\qquad
k_\tau=\frac{e^{-1.4\tau}}{6},\\
a_s=\min\left\{0.2,\frac{0.12}{\chi+0.4},
0.25-\psi_s\beta_1,0.25-\psi_s\beta_2\right\},\\
b_s=2k_\tau\psi_s\max\{\beta_1,\beta_2\}.
\end{gathered}
\]
If $a_s>b_s$, the linearization on the space of fixed total population
has exactly one neutral direction.
Its maximal Lyapunov exponent equals zero.
Every other exponent is at most
$-\min\{0.13,1,\eta_s\}<0$, where $\eta_s$ is the unique positive solution of
\begin{equation}
\eta_s=a_s-b_s e^{\eta_s\tau}.
\label{eq:lyapunov_decay_rate}
\end{equation}
\end{theorem}

\begin{proof}
Define $V(t)=1-U(t)$ and $P(t)=E(t)+H(t)+R(t)+D(t)$.
Population conservation gives $V(t)\ge0$.
The withdrawal equation gives
\[
\dot V(t)
=-0.2E-0.25R-0.25D
-\frac{0.3(M_1+M_2)}{\chi+M_1+M_2}H.
\]
The media equations imply $M_i(t)\ge0.2(1-e^{-t})$.
For $t\ge t_0=\log 2$, both media variables are at least $0.1$.
Hence
\[
\dot V(t)\le-c_\chi P(t),\qquad
c_\chi=\min\left\{0.2,0.25,\frac{0.06}{\chi+0.2}\right\}>0.
\]
Integration yields the quantitative bound
\begin{equation}
\int_{t_0}^{T}P(t)\,dt
\le\frac{V(t_0)-V(T)}{c_\chi}
\le\frac{1}{c_\chi}.
\label{eq:lyapunov_activity_bound}
\end{equation}
All states and their derivatives are bounded.
Thus $P$ is uniformly continuous.
A nonnegative, uniformly continuous function with a finite integral
converges to zero.
Therefore $E,H,R,D\to0$.

The media equations then give $M_1,M_2\to0.2$.
Let $S=S_u+S_e$.
Since $\dot S=-\mathcal I\le0$, we have $S\to s$ for some $s\in[0,1]$.
The susceptible equation takes the form
\[
\dot S_u=-c(t)S_u+0.25M_2S,\qquad
c(t)=0.4M_1+0.25M_2+\beta_1R+\beta_2D.
\]
Here $c(t)\to0.13$ and $0.25M_2S\to0.05s$.
Scalar linear stability gives $S_u\to5s/13$.
Population conservation gives the remaining limits in
\eqref{eq:closed_equilibrium_family}.
Convergence to a single equilibrium excludes sustained chaotic motion.

For the exponent claim, let $z=(v_E,v_H,v_R,v_D)^T$.
At $X^*(s)$, this subsystem is independent of the remaining perturbations.
It satisfies $\dot z=A_s z+B_s z(t-\tau)$, where
\[
A_s=
\begin{pmatrix}
-1.4&0&\psi_s\beta_1&\psi_s\beta_2\\
0&-h_*&0&0\\
0.8&\Gamma_R^*&-0.35&0.05\\
0.4&\Gamma_D^*&0.10&-0.30
\end{pmatrix},
\]
\[
B_s=k_\tau\psi_s
\begin{pmatrix}
0&0&-\beta_1&-\beta_2\\
0&0&\beta_1&\beta_2\\
0&0&0&0\\
0&0&0&0
\end{pmatrix}.
\]
Here $h_*=\Gamma_R^*+\Gamma_D^*+\Gamma_U^*$ and
$\Gamma_U^*=0.12/(\chi+0.4)$.
The logarithmic $1$-norm of $A_s$ is $-a_s$.
The induced $1$-norm of $B_s$ is $b_s$.
Consequently,
\[
D^+\|z(t)\|_1
\le-a_s\|z(t)\|_1
+b_s\sup_{u\in[t-\tau,t]}\|z(u)\|_1.
\]
When $a_s>b_s$, comparison with $C e^{-\eta_s t}$ gives
exponential decay with the rate in~\eqref{eq:lyapunov_decay_rate}.
The media perturbations have decay rate $1$ and receive only decaying input.
The susceptible block has eigenvalues $0$ and $-0.13$.
Its input from the active states and media decays exponentially.
The fixed population condition determines $v_U$.
Thus the only neutral direction is
\[
w=(5/13,8/13,0,0,0,0,-1,0,0)^T.
\]
Its constant history has exponent zero.
All remaining exponents satisfy the stated negative bound.
\end{proof}

We apply~\eqref{eq:lyapunov_growth_identity} to 12 neural parameter sets
and four independent DDE fits.
The 116 runs check mesh size, duration, perturbation direction,
initial history, and delay.
Each run extends to normalized time $500$.
We omit the first $100$ time units from the main estimate.
All estimates satisfy $|\widehat\lambda_{100,500}|<2.95\times10^{-8}$.
The estimates over $[300,500]$ are zero at the recorded numerical precision.

We also evaluate $a_s-b_s$ at the numerical terminal equilibria
from all 116 runs.
Its minimum is $0.0138776$.
The minimum value of $\eta_s$ is $0.0138124$.
Thus the sufficient decay condition holds at every evaluated equilibrium.
The resulting upper bound for all nonneutral exponents is
$-0.01381$ in inverse normalized time units.
This is a bound from the equilibrium linearization, not a computed
full Lyapunov spectrum or an interval certificate.
The near-zero dominant estimates agree with the neutral equilibrium direction.
They do not indicate sustained chaos.
The convergence theorem applies to this closed model with its stated coefficients.
It does not establish the dynamics of the open model or the observed events.
The supplement provides the numerical methods and convergence tables.
\endgroup
